\section*{Exercises}
\addcontentsline{toc}{section}{Exercises}

\subsection*{A. Theory}

\begin{exercise}
\exerciseid{4.A1}
Prove that the span of a list of vectors is a subspace.
\end{exercise}

\begin{exercise}
\exerciseid{4.A2}
Prove that if a list of vectors is linearly dependent, then at least one vector
in the list is a linear combination of the others.
\end{exercise}

\begin{exercise}
\exerciseid{4.A3}
Explain why a basis gives unique coordinates.
\end{exercise}

\subsection*{B. By Hand}

\begin{exercise}
\exerciseid{4.B1}
Do the vectors $(1,0,1)$, $(0,1,1)$, and $(1,1,2)$ form a basis of $\R^3$?
Explain without doing unnecessary arithmetic.
\end{exercise}

\begin{exercise}
\exerciseid{4.B2}
Let $b_1=(1,1)$ and $b_2=(1,-1)$. Find the coordinates of $v=(3,1)$ in the
basis $(b_1,b_2)$.
\end{exercise}

\begin{exercise}
\exerciseid{4.B3}
Find a basis for the span of
\[
  (1,2,0),\quad (2,4,0),\quad (0,1,1).
\]
\end{exercise}

\subsection*{C. Python / NumPy}

\begin{exercise}
\exerciseid{4.C1}
Put the vectors in Exercise 4.B3 into the columns of a matrix and use
\texttt{np.linalg.matrix\_rank} to check the dimension of their span.
\end{exercise}

\begin{exercise}
\exerciseid{4.C2}
Write a function \texttt{coords(B, v)} that returns the coordinate vector
$c$ solving $Bc=v$.
\end{exercise}

\subsection*{D. Challenge}

\begin{exercise}
\exerciseid{4.D1}
Give two different bases of $\mathcal P_2$ and compute the coordinates of
$p(x)=1+x+x^2$ in both bases.
\end{exercise}
